% ===================================================================
%  Chemie: Titrationskurve starke Säure / starke Base – BERECHNET
%  (einschließlich Autoprotolyse des Wassers), mit reinem TikZ.
%  Versuchsdaten oben ändern -> neue Kurve, Äquivalenzpunkt wandert mit.
%  Kompilieren: pdflatex vorlage_titration.tex
% ===================================================================
\documentclass[border=2mm]{standalone}
\usepackage[T1]{fontenc}
\usepackage{lmodern}
\usepackage[ngerman]{babel}
\usepackage{amsmath}
\usepackage[version=4]{mhchem}
\usepackage{tikz}

\definecolor{akzentblau}{HTML}{2563EB}
\definecolor{akzentrot}{HTML}{DC2626}

% ---- Versuchsdaten -------------------------------------------------
\def\cS{0.1}     % Konzentration der Salzsäure in mol/L
\def\VS{20}      % vorgelegtes Volumen der Säure in mL
\def\cB{0.1}     % Konzentration der Natronlauge in mol/L
% --------------------------------------------------------------------
% Überschuss d an H3O+ (positiv) bzw. OH- (negativ) in mol/L:
%   d = (cS*VS - cB*V) / (VS + V)
% Mit Autoprotolyse (Kw = 1e-14) gilt exakt  c(H3O+) = d/2 + sqrt(d^2/4 + Kw)
% und daraus  pH = 7 - arsinh(d / 2e-7) / ln 10.
% TikZ rechnet nur mit Zahlen bis etwa 16000 und 5 Nachkommastellen. Für
% |d| >= 1e-5 gilt arsinh(u) = sign(u) * ln(2|u|) bis auf weniger als 0,0004
% pH-Einheiten; ln(2|u|) = ln|d| + 7 ln 10 = ln|d| + 16.1181.
% Kleinere |d| liegen am Äquivalenzpunkt: dort ist pH = 7.
\pgfmathdeclarefunction{phwert}{1}{%
  \pgfmathparse{(\cS*\VS - \cB*#1)/(\VS + #1)}\let\d\pgfmathresult
  \pgfmathparse{abs(\d) < 0.00001 ? 7 : 7 - sign(\d)*(ln(max(abs(\d),0.00001)) + 16.1181)/ln(10)}}

\begin{document}
\begin{tikzpicture}[x=0.25cm, y=0.5cm, font=\small]
  \pgfmathsetmacro\Vaeq{\cS*\VS/\cB}          % Äquivalenzpunkt in mL
  \pgfmathsetmacro\Vmax{2*\Vaeq}
  % Gitter und Achsen
  \draw[dashed, black!20] (0,0) grid[xstep=5, ystep=2] (\Vmax,14);
  \draw[->] (0,0) -- (\Vmax+2,0);
  \node[below=18pt] at (\Vaeq,0) {zugegebenes Volumen $V(\ce{NaOH})$ in mL};
  \draw[->] (0,0) -- (0,15) node[above] {pH-Wert};
  \foreach \x in {0,5,...,\Vmax} \node[below] at (\x,0) {\pgfmathprintnumber[int detect]{\x}};
  \foreach \y in {2,4,...,14} \node[left] at (0,\y) {\y};
  % Kurve: 601 Punkte, darunter genau der Äquivalenzpunkt
  \draw[akzentblau, very thick, line join=round]
    plot[domain=0:\Vmax, samples=601] (\x,{phwert(\x)});
  \fill[akzentrot] (\Vaeq,7) circle (2.5pt) node[right=3pt] {Äquivalenzpunkt};
\end{tikzpicture}
\end{document}
